\documentclass[10pt,a4paper]{amsart}
\usepackage[margin=19mm]{geometry}
\usepackage{amsmath,amssymb,mathtools}
\usepackage[scr=rsfs,cal=euler]{mathalfa}
\usepackage{xfrac}
\usepackage[unicode,hidelinks,bookmarks=false]{hyperref}
\newcommand{\ie}{i.e.\ }
\newcommand{\cf}{cf.\ }
\newcommand{\resp}{respectively\ }
\newcommand{\Subsection}{\subsection}
\providecommand{\nobreakdash}{\nobreak\hbox{-}}
\setlength{\emergencystretch}{2em}
\multlinegap0pt
\usepackage{amssymb}
\usepackage{xcolor}
\usepackage{algorithm}\usepackage{algpseudocode}
\usepackage[shortlabels]{enumitem}
\newtheorem{lemma}{Lemma}
\newcommand{\E}{\EE}
\newcommand{\EE}{\mathbb{E}}
\newcommand{\R}{\RR}
\newcommand{\RR}{\mathbb{R}}
\newcommand{\abs}[1]{\lvert#1\rvert}
\newcommand{\bps}{{\rm BPS}}
\newcommand{\dd}{\ud}
\newcommand{\ip}[2]{\langle #1,#2\rangle}
\newcommand{\ud}{\,\mathrm{d}}
\title{Windowed thinning and query complexity for the bouncy particle and Zigzag samplers}
\author{Jianfeng Lu, Yinchen Luo}
\date{Source: arXiv:2607.28413v2 (CC BY 4.0)}
\begin{document}
\maketitle

\noindent\emph{Source and licence.} Windowed thinning and query complexity for the bouncy particle and Zigzag samplers. Version arXiv:2607.28413v2 (CC BY 4.0), distributed by arXiv under the Creative Commons Attribution 4.0 International licence, http://creativecommons.org/licenses/by/4.0/, as recorded in the arXiv licence metadata for this exact version and shown on the abstract page https://arxiv.org/abs/2607.28413v2. Full licence notice: \texttt{licenses/arxiv-2607.28413-CC-BY-4.0.txt}.\par\smallskip

\noindent\textit{Editorial scope.} Counted unit: the article's lemma lem:bps-count (source0.tex lines 1024--1039). The article's complete original proof of it is delivered with it (source0.tex lines 1041--1229, one \texttt{proof} environment of 189 lines). No mathematics is written, repaired or re-derived: the statement and the proof are the article's own lines, in the article's own notation, and no retained line is substituted or re-flowed.  Retained uncounted as the source-local context the proof needs: lines 185--188; lines 429--439; lines 527--529; lines 883--885; lines 893--896; lines 953--967; lines 1014--1019.  Nothing was omitted from any retained span, so the package carries no omission record.  No retained line contains a citation, so the wrapper carries no bibliography and no bibliography entry.  No retained line contains a figure environment, an image-inclusion command or drawing code, so no image is delivered and the package declares no asset dependency.  Editorial formatting only, in addition: an AMS \texttt{amsart} preamble that restates the article's own declarations and macros used by the retained lines, namely the environments \texttt{lemma} on separate counters, together with the standard packages those lines need.  The retained environments print in this document as Lemma 1 in order of appearance, whereas the article prints the same items with the numbering of its own full document.  The complete native source of the article as distributed in the arXiv e-print for this exact version is bundled unchanged as \texttt{sources/206416/source0.tex}.
\begin{equation}\label{eq:hessian-bounds}
        m I_d \preceq \nabla^2 U(x) \preceq L I_d,
        \qquad x\in\R^d.
\end{equation}

\begin{equation}\label{eq:bps-generator}
\mathcal L^{\bps} f(x,v)
=
v\cdot\nabla_x f(x,v)
+
\lambda^{\bps}(x,v)\bigl(f(x,R_x v)-f(x,v)\bigr)
+
\gamma^{\bps}\left(
  \int_{\R^d}f(x,w)\varphi(\dd w)-f(x,v)
\right).
\end{equation}

\begin{equation}\label{eq:cold-start}
        \rho_0= N(x_\star,L^{-1}I_d)\otimes N(0,I_d).
\end{equation}

\begin{equation}\label{eq:velocity-second-moment}
        \E\abs{V_t}^2=d.
\end{equation}

\begin{equation}\label{eq:position-second-moment}
        \bigl(\E\abs{X_t-x_\star}^2\bigr)^{1/2}
        \le\sqrt{d/L}+\sqrt d\,t.
\end{equation}

\begin{lemma}[Dynkin formula for observables of polynomial growth]
\label{lem:dynkin-polynomial}
Assume \eqref{eq:hessian-bounds} and the cold start
\eqref{eq:cold-start}, and consider either BPS or Zigzag, with generator
$\mathcal L$.  Let $f\in C^1(\mathbb R^{2d})$, and assume that both $f$ and $\mathcal Lf$ have polynomial growth; that is, for some $C<\infty$ and $p\ge0$,
\[
        \abs{f(x,v)}+\abs{\mathcal Lf(x,v)}
        \le C\bigl(1+\abs{x-x_\star}+\abs v\bigr)^p.
\]
Then, for every $T<\infty$,
\[
        \E f(X_T,V_T)-\E f(X_0,V_0)
        =\int_0^T\E[\mathcal Lf(X_t,V_t)]\dd t.
\]
\end{lemma}

\begin{equation}\label{eq:bps-motivating-jump}
        \psi(x,R_xv)-\psi(x,v)
        =-2\ip{v}{\nabla U(x)}
          \frac{\ip{x-x_\star}{\nabla U(x)}}
               {\abs{\nabla U(x)}^2}.
\end{equation}

\begin{lemma}[Bound on the BPS bounce count from a cold start]
\label{lem:bps-count}
Assume \eqref{eq:hessian-bounds}.
For the cold start \eqref{eq:cold-start} and every $T\ge0$,
\begin{equation}\label{eq:bps-square-rate-bound}
        \int_0^T\E\bigl[\lambda^{\bps}(X_t,V_t)^2\bigr]\dd t
        \;\le\; LdT+\frac{\gamma^{\bps} d}{4}+\frac{\sqrt L\,d}{2}.
\end{equation}
Hence, whenever $T\ge\max\{L^{-1/2},\,\gamma^{\bps}/(4L)\}$, the expected
number of bounces in $[0,T]$ satisfies
\begin{equation}\label{eq:bps-count-clean}
        B^{\bps}(T)
        =\int_0^T\E\bigl[\lambda^{\bps}(X_t,V_t)\bigr]\dd t
        \le 2\sqrt{Ld}\,T .
\end{equation}
\end{lemma}

\begin{proof}
Let us introduce the shorthand
\[
        r(x):=
        \begin{cases}
        \displaystyle
        \frac{\ip{x-x_\star}{\nabla U(x)}}{\abs{\nabla U(x)}^2},
                &\nabla U(x)\ne0,\\[2mm]
        1/L,    &\nabla U(x)=0.
        \end{cases}
\]
With this notation, \eqref{eq:bps-motivating-jump} can be written as
\[
        \psi(x,R_xv)-\psi(x,v)
        =-2\ip{v}{\nabla U(x)}r(x).
\]
Strong convexity implies that $\nabla U(x)=0$ only at $x_\star$, where the
bounce rate vanishes; the assigned value at $x_\star$ is therefore not important.  

We claim that
\begin{equation}\label{eq:bps-psi-generator}
        \mathcal L^{\bps}\psi(x,v)
        =\abs v^2-2\lambda^{\bps}(x,v)^2r(x)
          -\gamma^{\bps}\psi(x,v).
\end{equation}
Indeed, since $\nabla_x\psi(x,v)=v$, the transport contribution in
\eqref{eq:bps-generator} is
\[
        v\cdot\nabla_x\psi(x,v)=\abs v^2.
\]
The bounce term follows from the calculation above. Finally, since $\varphi$ is centered,
\[
        \int_{\R^d}\psi(x,w)\varphi(\dd w)
        =\ip{x-x_\star}{\int_{\R^d}w\varphi(\dd w)}
        =0,
\]
so the refreshment contribution is
$-\gamma^{\bps}\psi(x,v)$.  Combining the three contributions gives
\eqref{eq:bps-psi-generator}.

The identity \eqref{eq:bps-psi-generator} involves
the squared bounce rate weighted by $r$.  To recover the unweighted rate, we
need a uniform lower bound on this factor.  We claim that
\[
        \frac1L\le r(x)\le\frac1m,
        \qquad x\in\R^d.
\]
The claim holds at $x_\star$ by the definition of $r$.  Suppose now that
$x\ne x_\star$.  Since $\nabla U(x_\star)=0$, the fundamental theorem of
calculus gives
\[
        \nabla U(x)=A_x(x-x_\star),
        \qquad
        A_x:=\int_0^1
        \nabla^2U\bigl(x_\star+s(x-x_\star)\bigr)\dd s.
\]
$A_x$ is symmetric and satisfies
$mI_d\preceq A_x\preceq L I_d$ by the bound on Hessian of $U$.  Its spectrum is therefore contained in
$[m,L]$, and hence $mA_x\preceq A_x^2\preceq LA_x$. Thus we arrived at the desired bound for $r$ combined with the observation that 
\[
\begin{aligned}
        r(x)
        =\frac{\ip{x-x_\star}{\nabla U(x)}}
                {\abs{\nabla U(x)}^2}
        =\frac{\ip{x-x_\star}{A_x(x-x_\star)}}
                {\ip{x-x_\star}{A_x^2(x-x_\star)}}.
\end{aligned}
\]

We may now apply Dynkin's formula.  The bounds
$\lambda^{\bps}(x,v)\le L\abs{x-x_\star}\abs v$ and
$r(x)\le1/m$ show from \eqref{eq:bps-psi-generator} that
$\abs\psi+\abs{\mathcal L^{\bps}\psi}$ has polynomial growth.  Therefore
Lemma~\ref{lem:dynkin-polynomial} applies and yields
\[
\begin{aligned}
        \E\psi(X_T,V_T)-\E\psi(X_0,V_0)
        & =\int_0^T
          \E\!\left[\mathcal L^{\bps}\psi(X_t,V_t)\right]\dd t \\
        &=\int_0^T\E\abs{V_t}^2\dd t
          -2\int_0^T\E\!\left[
            \lambda^{\bps}(X_t,V_t)^2r(X_t)
          \right]\dd t\\
        &\quad-\gamma^{\bps}
          \int_0^T\E\psi(X_t,V_t)\dd t,
\end{aligned}
\]
where the second equality uses  \eqref{eq:bps-psi-generator}.
By \eqref{eq:velocity-second-moment}, the first term on the right is $dT$.
Rearranging the preceding identity therefore gives
\begin{equation}\label{eq:dynkin}
\begin{aligned}
        2\int_0^T\E\!\left[
          \lambda^{\bps}(X_t,V_t)^2r(X_t)
        \right]\dd t
        &=dT-\gamma^{\bps}\int_0^T\E\psi(X_t,V_t)\dd t\\
        &\quad-\E\psi(X_T,V_T)+\E\psi(X_0,V_0).
\end{aligned}
\end{equation}
We next eliminate the time integral of $\psi$.  The position is continuous,
and $\dot X_t=V_t$ between events.  Therefore, for almost every $t$,
\[
        \frac{\dd}{\dd t}\abs{X_t-x_\star}^2
        =2\ip{X_t-x_\star}{V_t}
        =2\psi(X_t,V_t).
\]
Integration over $[0,T]$ and taking expectation therefore gives 
\[
        \int_0^T\E\psi(X_t,V_t)\dd t
        =\frac12\left(
          \E\abs{X_T-x_\star}^2-\E\abs{X_0-x_\star}^2
        \right).
\]
Consequently, the time-integral term in the rearranged Dynkin identity is
\[
        -\gamma^{\bps}\int_0^T\E\psi(X_t,V_t)\dd t
        =-\frac{\gamma^{\bps}}2
          \left(\E\abs{X_T-x_\star}^2-\E\abs{X_0-x_\star}^2\right).
\]
We now estimate this term together with the two endpoint terms on the right hand side of 
identity \eqref{eq:dynkin}. Under the cold start, $X_0-x_\star$ and $V_0$ are independent and centered, and hence
\[
        \E\psi(X_0,V_0)
        =\E\ip{X_0-x_\star}{V_0}=0,
        \qquad
        \E\abs{X_0-x_\star}^2=\frac dL.
\]
Since $\E\abs{X_T-x_\star}^2\ge0$, it follows that
\[
        -\gamma^{\bps}\int_0^T\E\psi(X_t,V_t)\dd t
        \le\frac{\gamma^{\bps}d}{2L}.
\]
For the remaining endpoint term, Cauchy--Schwarz,
\eqref{eq:velocity-second-moment}, and
\eqref{eq:position-second-moment} give
\[
\begin{aligned}
        -\E\ip{X_T-x_\star}{V_T}
        &\le
        \left|\E\ip{X_T-x_\star}{V_T}\right|\\
        &\le
        \sqrt{\E\abs{X_T-x_\star}^2\,\E\abs{V_T}^2}
        \le\left(\sqrt{\frac dL}+\sqrt d\,T\right)\sqrt d =\frac d{\sqrt L}+dT.
\end{aligned}
\]
Inserting these three bounds into \eqref{eq:dynkin} yields
\[
\begin{aligned}
        2\int_0^T\E\!\left[
          \lambda^{\bps}(X_t,V_t)^2r(X_t)
        \right]\dd t
        &\le 2dT+\frac{\gamma^{\bps}d}{2L}+\frac d{\sqrt L},
\end{aligned}
\]
and the lower bound $r\ge1/L$ then gives
\[
        \int_0^T
        \E\bigl[\lambda^{\bps}(X_t,V_t)^2\bigr]\dd t
        \le LdT+\frac{\gamma^{\bps}d}{4}+\frac{\sqrt L\,d}{2},
\]
which proves \eqref{eq:bps-square-rate-bound}.

It remains to deduce the bounce-count estimate. Cauchy--Schwarz gives
\[
        B^{\bps}(T)
        =\int_0^T\E\lambda^{\bps}(X_t,V_t)\dd t
        \le\left(
          T\int_0^T\E[\lambda^{\bps}(X_t,V_t)^2]\dd t
        \right)^{1/2}.
\]
If $T\ge L^{-1/2}$ and $T\ge\gamma^{\bps}/(4L)$, then
\[
        \frac{\sqrt L\,d}{2}\le\frac12LdT,
        \qquad
        \frac{\gamma^{\bps}d}{4}\le LdT.
\]
Thus \eqref{eq:bps-square-rate-bound} implies
\[
        \int_0^T\E[\lambda^{\bps}(X_t,V_t)^2]\dd t
        \le\frac52LdT.
\]
Substitution into the Cauchy--Schwarz bound gives
\[
        B^{\bps}(T)
        \le\sqrt{\frac52}\,\sqrt{Ld}\,T
        \le2\sqrt{Ld}\,T,
\]
which proves \eqref{eq:bps-count-clean}.
\end{proof}

\end{document}
